\begin{proof}
Let $u_i=\log(p_i/q_i)$ and, with $I\sim p$, put
$W=\max_i u_i-u_I\in[0,e]$. Since $\mathbb E_p e^{-u_I}=1$,
\[
 D_{\rm KL}(p\|q)=\log\mathbb E e^W-\mathbb EW.
\]
For a fixed mean $\mu\in[0,e]$, convexity gives
$\mathbb E e^W\le1+(\mu/e)(\exp(e)-1)$, with equality for an
endpoint-supported law. Maximizing the resulting scalar expression
over $\mu$ gives $\mu_*=1-e/(\exp(e)-1)$ and the stated $\Psi(e)$.
To see attainability, give $W\in\{0,e\}$ mean $\mu_*$, call its
probabilities $p_i$, set $b=\log\mathbb E e^W$, and define
$q_i=p_i\exp(W_i-b)$. These $q_i$ sum to one, are positive, and
$\log(p_i/q_i)=b-W_i$ has oscillation $e$.

For the quadratic bound let $K(t)=\log\mathbb E e^{tW}$.
Its second derivative is the variance of $W$ under an exponential
tilt. Any variable in $[0,e]$ has variance at most $e^2/4$, since
$\operatorname{Var}(W)\le\mathbb E(W-e/2)^2\le e^2/4$.
Integrating $K''$ with weight $1-t$ on $[0,1]$ yields
$K(1)-K(0)-K'(0)\le e^2/8$.
Expanding the displayed elementary formula gives the asserted series;
using $e/(\exp(e)-1)\sim e\exp(-e)$ gives the large-$e$ expression.
\end{proof}
